\section{Introduction}
\label{sec:intro}

The proliferation of task-specific, fine-tuned models derived from shared pre-trained initializations has catalyzed intense interest in weight-space model merging and adaptive ensembling~\cite{ainsworth2022git, wortsman2022model, yadav23ties, yu23dare}. While joint multi-task training is often computationally prohibitive, model merging offers a cheap, modular, and post-hoc mechanism to combine independent expert networks by interpolating their parameters in weight space. A major focus is \emph{layer-wise adaptive merging}, where ensembling coefficients $\alpha_k(l)$ are tuned across task experts $k$ and network layers $l$ to construct a dynamic, unified architecture~\cite{choshen22fusing, chatterjee2024robustness}. When optimized on small calibration datasets (e.g., in few-shot or test-time adaptation), this high-dimensional parameter space is highly prone to transductive overfitting, leading to poor generalization on the true multi-task distribution.

To prevent this overfitting, prior methods have attempted to restrict the ensembling trajectory across depth. For instance, Rademacher-Bounded Polynomial Merging (RBPM) projects layer-wise coefficients onto a low-degree polynomial subspace, applying a learning-theoretic penalty based on the polynomial degree $d$~\cite{chatterjee2024robustness}. However, low-degree polynomial trajectories (such as RBPM, $d=2$) suffer from severe boundary runaway during optimization. Due to the rigid, global parabolic shape constraints of quadratic curves on $z \in [0, 1]$, fitting intermediate representation layers forces the coefficients at the boundary domains (i.e., the first and last layers of deep backbones) to extreme values. Because the first layers govern crucial low-level feature extraction and the final layers project representations into task logits, boundary runaway catastrophically degrades classification accuracy.

In this paper, we resolve these theoretical and practical limitations by introducing \textbf{Rademacher-Bounded Fourier Trajectory Merging (RB-FTM)} and its non-periodic counterpart, \textbf{Rademacher-Bounded Discrete Cosine Trajectory Merging (RB-DCTM)}. Drawing from Fourier analysis and statistical learning theory, we parameterize the ensembling trajectory across depth as a continuous spectral series composed of a small number of low-frequency harmonic sinusoids. Since sinusoidal basis functions are inherently bounded in $[-1, 1]$ and represent periodic harmonic waves, their values are highly stable, completely mitigating boundary runaway. To isolate ensembling dynamics from confounding hardware and optimization noise, we conduct our analysis inside a mathematically controlled, synthetic \textbf{Analytical Coordinate Sandbox (ACS)}, which simulates representation propagation through deep layer coordinates. Our evaluations show that while the parameter-free, zero-tuning \textbf{Static Uniform} ensembling baseline remains a highly robust consensus upper bound in perfectly coordinate-aligned spaces (due to perfect representation alignment), our proposed spectral trajectory models serve as exceptionally robust regularized optimizers when adaptation of ensembling weights is required. Specifically, RB-FTM and RB-DCTM significantly outperform all other tuned and adaptive baselines (such as unconstrained, globally-scaled, and polynomial baselines) across both simulated convolutional and transformer backbone architectures, while maintaining smooth, stable parameter trajectories.

Our key contributions are summarized as follows:
\begin{itemize}
    \item \textbf{Trigonometric Representation of Merging Trajectories}: We project discrete layer-wise ensembling coefficients onto a low-frequency continuous Fourier subspace across depth, acting as a natural low-pass filter that prunes high-frequency noise and representation misalignment.
    \item \textbf{Learning-Theoretic Complexity Guarantees}: We prove that the empirical Rademacher complexity of our proposed Fourier trajectory class is strictly bounded by $C_0 \sqrt{2 \ln(4F+2) / L}$, where $F$ is the spectral cutoff and $L$ is the depth. This bound mathematically limits the structural complexity of the trajectory class across layers independent of parameter count, acting as a powerful coefficient-space regularizer.
    \item \textbf{Analytical Spectral Lasso Regularization}: We formulate an optimization objective that directly couples our learning-theoretic bound to a practical $L_1$ penalty on the Fourier coefficients. This acts as a physical regularizer during few-shot calibration, compressing the active spectral capacity.
    \item \textbf{Runge's Phenomenon Mitigation and Empirical Validation}: We mathematically and empirically demonstrate that RB-FTM and RB-DCTM eliminate the boundary oscillations of polynomial trajectories. Evaluations in our Analytical Coordinate Sandbox across Deep12LayerCNN and CLIP ViT-B/16 on a four-task stream show that our proposed spectral methods achieve substantial accuracy gains over all other tuned baselines while generating highly interpretable, smooth trajectories. We also provide a small-scale real-world validation experiment on actual Vision Transformer parameters.
\end{itemize}